\section{从投票自治到分层治理}

早期 Reddit 的原则接近“平台不删除，用户用投票决定可见性”。在规模较小、恶意参与有限时，这种自治降低了中心化编辑权力，也鼓励社区形成自己的文化。规模扩大后，投票无法处理跨社区骚扰、非法内容、隐私泄露和专门以伤害为目标的群体。原因并不是用户突然不可信，而是伤害具有外部性：受影响者可能不在投票社区内，少数协调者也能操纵局部多数。

\subsection{治理为什么必须分层}

平台级规则负责最低安全、合法性和跨社区边界，subreddit 版主负责本地主题与规范，用户投票负责日常排序和反馈。三层分别拥有不同信息：平台能看见跨站模式和法律风险，版主理解语境与社区历史，普通用户最接近内容价值。让任一层独占全部权力都会失真，治理系统的目标是让问题在拥有足够信息且影响范围匹配的层级处理。

\begin{figure}[H]
\centering
\includegraphics[width=0.92\textwidth]{images/03-governance-stack.png}
\caption{社区治理栈：全站政策、平台执法、subreddit 规则与用户参与分工。}
\figsource{依据访谈字幕 08:30--15:00 重绘，`images/03-governance-stack.png`。}
\end{figure}

\begin{knowledgebox}{读图：层级越高，规则应越少但边界越硬}
Site-wide policy 面向跨站伤害，必须相对稳定、可解释并受法律与公共责任约束。Subreddit rules 可以更细，允许不同社区对发帖格式、话题和互动风格作出选择。Voting 是高频弱信号，适合排序，不适合判定所有安全问题。治理失灵常发生在层级错配：让投票处理人肉搜索，或让平台总部决定每个小社区的语气。
\end{knowledgebox}

\begin{importantbox}{首次使用术语：moderation 与 doxing}
Moderation 是内容与社区治理，包括制定规则、审核举报、限制账号、处理申诉和支持版主。Doxing 指未经同意公开个人身份、住址、联系方式等敏感信息，可能把线上冲突转化为现实伤害。它不能仅靠受欢迎程度决定，需要全站级限制和快速响应。
\end{importantbox}

\subsection{``Specifically vague''：规则为何不能写成穷举表}

恶意参与者会寻找字面漏洞。如果规则只列出固定词语或行为，攻击者可以改变拼写、拆分行动或组织他人完成同一伤害。规则太宽又会让普通用户无法预测边界，执法团队也容易不一致。Huffman 用 ``specifically vague'' 描述工作区间：清楚说明要避免的伤害与原则，同时保留根据语境判断等价规避行为的空间。

\begin{figure}[H]
\centering
\includegraphics[width=0.92\textwidth]{images/04-policy-specificity.png}
\caption{政策设计需要在可预测性与抵抗字面规避之间寻找工作区间。}
\figsource{依据访谈字幕 13:40--16:00 重绘，`images/04-policy-specificity.png`。}
\end{figure}

\begin{knowledgebox}{读图：模糊空间必须由程序约束}
中间区域不是允许执法者随意决定。它需要示例、先例、跨团队审查、申诉、透明报告和定期修订来约束。越依赖语境判断，越要记录为什么相似案件得到不同处理。政策文本负责表达目标，程序正义负责让目标不被个人偏见占据。
\end{knowledgebox}

\begin{warningbox}{课堂提示：每个词都来自一次代价}
讲者强调，成熟内容政策中的词句往往是多年事故、争议和边界案例压缩后的结果。读者看到十条规则时，不应只评估文案是否优雅，还应追问相应检测、人工审查、通知、申诉和恢复工具是否存在。没有执行系统的规则只会把责任推给一线团队。
\end{warningbox}

\subsection{执法从二元开关变成梯子}

早期工具接近永久封禁或不处理。二元选择会让团队过度执法，也会因为害怕不可逆损失而完全不行动。更成熟的系统把解释、警告、短期限制、subreddit timeout、临时停权和永久移除组成梯子。干预强度根据伤害、重复次数、合作意愿和恢复可能性上升，并为误判提供回退路径。

\begin{figure}[H]
\centering
\includegraphics[width=0.94\textwidth]{images/05-enforcement-ladder.png}
\caption{分级执法：从解释与警告逐步升级到临时限制和永久移除。}
\figsource{依据访谈字幕 16:00--20:00 重绘，`images/05-enforcement-ladder.png`。}
\end{figure}

\begin{importantbox}{读图：工具精度决定治理精度}
如果系统只有永久 ban，政策团队即使理解差异也无法表达差异。Timeout 让平台暂停一个失控社区、与版主沟通并保留恢复可能；递增停权让偶发违规者获得纠正机会，同时为持续恶意行为积累证据。治理质量因此是政策、数据和产品工具的共同结果。
\end{importantbox}

\begin{table}[H]
\centering
\begin{tabular}{L{0.18\textwidth}L{0.27\textwidth}L{0.43\textwidth}}
\toprule
干预 & 适用场景 & 必须配套的程序 \\
\midrule
解释/提醒 & 初次、低危、规则不清 & 指向具体规则与内容 \\
警告 & 已确认违规但可纠正 & 留存记录并说明再次违规后果 \\
临时限制 & 重复违规或短期风险 & 明确期限、范围和恢复条件 \\
社区 timeout & 版主失去控制或冲突升级 & 与版主沟通、保护历史内容、复开审查 \\
永久移除 & 严重伤害、恶意规避或持续失败 & 高级复核、证据保存、申诉与透明报告 \\
\bottomrule
\end{tabular}
\caption{执法梯子把伤害、可逆性和程序保障放在同一张表中。}
\end{table}

\subsection{本章小结}

投票适合排序，不足以承担全部安全与合法性责任。Reddit 的治理演进把全站政策、平台执法、社区规则和用户参与分层，并用 ``specifically vague'' 与分级工具处理恶意规避和边界案例。下一章讨论另一种边界：公开可读内容是否意味着任何主体都能无限量抓取和商业利用。

